% - Want to recover multi-material physical properties through a non-differentiable fracture simulator
% - Talk about the anisotropy? - Maybe in the background or problem formulation or method

% Simulating realistic food fracture requires specifying physical material parameters that are difficult to determine from first principles. We formulate the inverse problem: given a description of desired fracture behavior, automatically discover the material parameters that produce it. Our test case is orange peeling simulated with AnisoMPM, a continuum damage mechanics simulator whose fracture mechanics are non-differentiable, ruling out gradient-based inverse methods. We propose a learned inverse mapping that takes a target behavioral description and directly outputs material parameters, trained as a goal-conditioned reinforcement learning policy operating in the latent space of a normalizing flow. The flow compresses a 9-dimensional physical parameter space to 4 dimensions while preserving an invertible structure that avoids the optimization pathologies of variational autoencoders. To train and evaluate our approach, we generate 2,000 forward simulations, extract aggregate behavioral features, and train a neural surrogate that accelerates evaluation by a factor of $1.5 \times 10^6$. We compare our learned inverse mapping against CMA-ES operating in both the raw parameter space and the flow latent space, and find that the policy produces competitive parameter estimates in a single forward pass for targets where CMA-ES requires thousands of surrogate evaluations. We also identify a significant gap between surrogate-predicted and simulator-validated performance, and show that searching in the flow latent space substantially outperforms the equivalent search in a VAE latent space for material parameter recovery.

\begin{abstract}
Realistic visual simulation of food manipulation requires accurate material parameters, yet these are difficult to measure directly and vary across the heterogeneous regions of a single food item.
We address the inverse problem of estimating material parameters from a target description of fracture behavior in a non-differentiable continuum damage mechanics simulator.
Using orange peeling as a test case, we train a neural surrogate on 2,000 forward simulations and compare Covariance Matrix Adaptation Evolution Strategy (CMA-ES, a gradient-free evolutionary optimizer) with Proximal Policy Optimization (PPO, a reinforcement learning algorithm) across the original 9-dimensional parameter space and two learned 4-dimensional latent representations.
Since different oranges have different material properties, a practical inverse system must handle arbitrary targets without retraining. We train a goal-conditioned PPO policy that learns a general inverse mapping: given any target description of peeling behavior, the policy produces a material parameter estimate in a single forward pass (8 surrogate evaluations, approximately 10\,ms).
Operating in a normalizing flow latent space with a shared surrogate evaluator, the goal-conditioned policy achieves 0.642 actual recovery when validated through the simulator, outperforming the original parameter space by 23\%.
A warm-start extension that initializes CMA-ES refinement from the policy's output further improves recovery to 0.828 with 540 evaluations.
These findings provide a practical framework for inverse food physics and lay groundwork for vision-driven material identification from video observations of food manipulation.


\end{abstract}